\documentclass[conference]{IEEEtran}
\usepackage{cite,amsmath,amssymb,graphicx,booktabs,url,hyperref,xcolor}
\newcommand{\todo}[1]{\textcolor{red}{[TODO: #1]}}

\begin{document}

\title{Explainable Hybrid Intrusion Detection:\\Combining Declarative Signatures with Per-Feature Anomaly Attribution}

\author{\IEEEauthorblockN{\todo{Author}}
\IEEEauthorblockA{School of Electronics and Computer Science\\University of Southampton\\\todo{email}}}

\maketitle

\begin{abstract}
\todo{150--200 words. Problem: anomaly-based NIDS reach high recall but produce alerts analysts cannot act on; signature systems are precise but miss novel attacks. Approach: a hybrid pipeline whose anomaly path reports which flow feature triggered the alert, fused with declarative rules through a shared alert manager. Evaluation: CIC-IDS2017 and UNSW-NB15, per-attack precision/recall, false-positive rate at matched recall, alert-volume reduction from de-duplication. Headline result: \todo{numbers}.}
\end{abstract}

\begin{IEEEkeywords}
intrusion detection, anomaly detection, network security, explainability, flow features
\end{IEEEkeywords}

\section{Introduction}
\todo{Motivation: alert fatigue; the precision/recall gap between signature and anomaly NIDS \cite{sommer2010outside}. State the research question (see \texttt{docs/RESEARCH\_PLAN.md}). List contributions as 3--4 bullets.}

\section{Background and Related Work}
\subsection{Signature-based detection}
\todo{Snort \cite{roesch1999snort}, Suricata; strengths and the novel-attack blind spot.}
\subsection{Anomaly-based detection}
\todo{Statistical baselines, Isolation Forest \cite{liu2008isolation}, deep models; the operational-deployment critique \cite{sommer2010outside}.}
\subsection{Hybrid and explainable NIDS}
\todo{Prior hybrids; explainability for IDS \cite{arp2022dos}; where this work sits.}
\subsection{Benchmark datasets}
\todo{CIC-IDS2017 \cite{sharafaldin2018toward}, UNSW-NB15 \cite{moustafa2015unsw}, NSL-KDD; known labelling issues \cite{engelen2021troubleshooting}.}

\section{Threat Model and Design Goals}
\todo{Passive sensor on one link; external attacker; sensor uncompromised. Goals G1--G4 from DESIGN.md.}

\section{System Design}
\subsection{Flow assembly and features}
\todo{5-tuple bidirectional flows, idle timeout $\tau$, feature vector $\mathbf{x} \in \mathbb{R}^6$ (Table~\ref{tab:features}).}
\begin{table}[t]
\centering\caption{Per-flow features.}\label{tab:features}
\begin{tabular}{ll}\toprule
Feature & Definition \\\midrule
duration & $t_{\text{last}} - t_{\text{first}}$ \\
packets & packet count \\
bytes & total bytes \\
mean\_len & bytes / packets \\
syn\_ratio & SYN-only packets / packets \\
pkts\_per\_sec & packets / duration \\\bottomrule
\end{tabular}
\end{table}
\subsection{Signature path}
\todo{Declarative rule language; matching semantics.}
\subsection{Anomaly path with feature attribution}
Each flow is scored by $z_j = |(\log(1+x_j) - \mu_j)/\sigma_j|$ and an alert is raised when $\max_j z_j \ge \theta$, reporting the arg-max feature $j^\star$ as the explanation. \todo{Justify log transform; contrast with Isolation Forest, which gives a score but no attribution.}
\subsection{Alert fusion and de-duplication}
\todo{Dedupe key, window $w$, severity ordering.}

\section{Experimental Setup}
\todo{Datasets and splits; attack classes; baselines (signature-only, anomaly-only, Isolation Forest); metrics (precision, recall, F1, FPR, alerts per true incident); hyperparameters $\tau,\theta,w$; hardware.}

\section{Results}
% Regenerated by: python -m ids.evaluation.run --dataset cicids2017
\todo{Run the evaluation harness to populate this table.}

\todo{Per-attack table; precision--recall curves over $\theta$ (Fig.~\ref{fig:pr}); alert-volume reduction; ablation of the log transform; attribution agreement with ground truth attack type.}
\begin{figure}[t]\centering
\todo{\includegraphics[width=\linewidth]{figures/pr_curve.pdf}}
\caption{Precision--recall over anomaly threshold $\theta$.}\label{fig:pr}
\end{figure}

\section{Discussion}
\todo{Why attribution matters operationally; where the hybrid fails (low-rate attacks, encrypted payloads); dataset artefacts.}

\section{Threats to Validity}
\todo{Internal: label noise in CIC-IDS2017; feature set small. External: single-link, offline evaluation; synthetic traffic in unit tests only. Construct: FPR per flow vs per incident.}

\section{Conclusion and Future Work}
\todo{Summarise; live capture, per-service baselines, learned rule mining.}

\section*{Reproducibility}
Code, rules, configs and the evaluation harness are in the \texttt{ids/} directory of the project repository. \texttt{python -m ids.evaluation.run} regenerates every table in this paper from the dataset CSVs.

\bibliographystyle{IEEEtran}
\bibliography{references}
\end{document}
